\section*{Note de synthèse}
\addcontentsline{toc}{section}{Note de synthèse}

\textbf{Cadre.} Ce travail a été mené pendant le stage de fin d'études réalisé chez
\acompleter{nom de l'entreprise}, au sein de \acompleter{équipe}, du \acompleter{dates}. Il porte sur
l'évaluation d'options sur obligations d'État françaises, les OAT. Une option d'achat (call) donne le droit
d'acheter l'obligation à une date future à un prix fixé aujourd'hui ; une option de vente (put), celui de la
vendre. Le contrat étudié est celui qui se traite en banque. L'option est garantie par un dépôt de collatéral
rémunéré au taux interbancaire au jour le jour €STR, et son prix à terme se déduit du coût de financement du titre
sur le marché de la pension livrée (repo).

\textbf{Problématique.} Le prix d'une OAT dépend de deux sources de risque : le niveau des taux sans risque, et
l'écart (spread) entre le rendement de l'OAT et ces taux, qui rémunère le risque de crédit de l'État, la liquidité
du titre et des effets d'offre. Les modèles de taux usuels ne retiennent que la première source. Or il n'existe pas
de marché d'options sur OAT d'où l'on pourrait tirer la volatilité du spread. La question posée est double :
faut-il modéliser le spread comme un risque à part pour évaluer une option sur OAT, et les données disponibles
permettent-elles d'en mesurer les caractéristiques ?

\textbf{Démarche.} Le modèle retenu, issu de la littérature (Brigo et Mercurio, 2006 ; Russo, Giacometti et
Fabozzi), décrit le taux de l'OAT comme la somme de deux facteurs aléatoires gaussiens : le taux sans risque et le
spread. Le facteur taux est calibré sur les prix des swaptions, options sur taux d'intérêt très liquides. Pour le
facteur spread, on montre d'abord que la méthode proposée dans l'article de référence, fondée sur la seule courbe
des taux OAT, ne peut rien mesurer. Les paramètres du spread sont donc estimés sur six ans de données
quotidiennes (2021-2026), par deux méthodes économétriques enseignées à l'ENSAE : la méthode des moments
généralisée et le maximum de vraisemblance avec filtre de Kalman. Le prix de l'option est calculé exactement, et
ce calcul est validé par simulation de Monte-Carlo. Les données proviennent de Bloomberg (8 septembre 2026), et
l'option porte sur l'OAT 4,10~\% 2046, pour des échéances de un à dix ans.

\textbf{Résultats.} Quatre résultats se dégagent.
\begin{itemize}
\item La volatilité du spread est d'environ 34 points de base par an, sa corrélation avec les taux de 0,13 ;
elle alterne entre une période calme (23 points de base) et des épisodes agités (51 points de base), comme
en 2022 ou lors des crises politiques de 2024, mais ces épisodes sont trop courts pour modifier le prix.
\item Prendre en compte le spread renchérit l'option de 17~\%, à toutes les échéances.
\item Cet effet se résume, pour une option sur une seule OAT, à une volatilité augmentée : un modèle à un seul
facteur, dont la volatilité est relevée d'environ 17~\%, reproduit les prix à 0,05~\% près.
\item La formule approchée de l'article de référence est précise pour les options à la monnaie, mais commet des
erreurs de 3 à 4~\% sur les options éloignées de la monnaie : le prix exact est préférable.
\end{itemize}
Enfin, pour les échéances longues, le coût de financement en repo pèse autant sur le prix que le spread lui-même.

\textbf{Limites et suites.} Faute de cotations, la courbe de financement en repo utilisée est fictive ; une courbe
de marché la remplacerait sans modifier les calculs. Les tests statistiques montrent qu'un seul facteur par courbe
ne décrit pas toutes les déformations du spread, et la corrélation entre spread et taux est instable dans le
temps. Ces points ouvrent la voie à des modèles plus riches, utiles surtout pour des produits sensibles à l'écart
entre plusieurs obligations.
